% Bagian: first-order-logic
% Bab: proof-systems
% Subbagian: tableaux

\documentclass[../../../include/open-logic-section]{subfiles}

\begin{document}

\iftag{FOL}
      {\olfileid[id]{fol}{prf}{tab}}
      {\olfileid[id]{pl}{prf}{tab}}

\olsection{\usetoken{P}{tableau}}

Sementara banyak sistem !!{derivation} bekerja dengan susunan
!!{sentence}s, !!{tableau}s bekerja dengan !!{signed formula}s.
!!^a{signed formula} adalah pasangan yang terdiri atas tanda nilai
kebenaran ($\True$ atau $\False$) dan !!a{sentence}
\[
\sFmla{\True}{!A} \text{ atau } \sFmla{\False}{!A}.
\]
!!^a{tableau} terdiri atas !!{signed formula}s yang disusun dalam
pohon bercabang ke bawah. Tableau itu berawal dengan sejumlah
\emph{asumsi} dan dilanjutkan dengan !!{signed formula}s yang diperoleh
dari salah satu !!{signed formula}s di atasnya melalui penerapan salah
satu aturan inferensi. Setiap aturan memungkinkan kita menambahkan
satu atau lebih !!{signed formula}s pada ujung suatu cabang, atau dua
!!{signed formula}s secara berdampingan---dalam hal ini sebuah cabang
terbelah menjadi dua, dengan dua !!{signed
  formula}s yang ditambahkan membentuk ujung kedua cabang tersebut.

Aturan penghubung yang diterapkan pada !!{signed formula} kompleks
menghasilkan penambahan !!{signed formula}s yang merupakan
sub-!!{formula}s langsung; aturan kuantor dapat menghasilkan instans
substitusi dari suatu subformula. Aturan-aturan itu muncul berpasangan,
satu aturan untuk masing-masing dari kedua tanda. Sebagai contoh, aturan
$\TRule{\True}{\land}$ diterapkan pada $\sFmla{\True}{!A \land !B}$,
dan memungkinkan penambahan kedua !!{signed formula}s, yaitu
$\sFmla{\True}{!A}$ dan~$\sFmla{\True}{!B}$, pada ujung setiap cabang
yang memuat $\sFmla{\True}{!A \land !B}$, sedangkan aturan
$\TRule{\False}{\land}$ memungkinkan suatu cabang dibelah dengan
menambahkan $\sFmla{\False}{!A}$ dan $\sFmla{\False}{!B}$ secara
berdampingan. !!^a{tableau} tertutup jika setiap cabangnya memuat
pasangan !!{signed formula}s yang cocok, yakni
$\sFmla{\True}{!A}$ dan $\sFmla{\False}{!A}$.

Relasi $\Proves$ berdasarkan !!{tableau}s didefinisikan sebagai
berikut: $\Gamma \Proves !A$ jika dan hanya jika terdapat suatu
himpunan berhingga~$\Gamma_0 = \{!B_1, \dots, !B_n\} \subseteq \Gamma$
sedemikian sehingga terdapat !!{tableau} tertutup untuk asumsi-asumsi
\[
\{\sFmla{\False}{!A}, \sFmla{\True}{!B_1}, \dots, \sFmla{\True}{!B_n}\}
\]
Sebagai contoh, berikut adalah suatu !!{tableau} tertutup yang menunjukkan
bahwa $\Proves (!A \land !B) \lif !A$:
\begin{oltableau}
  [\sFmla{\False}{(\formula{A} \land \formula{B}) \lif \formula{A}}, just = \TAss
    [\sFmla{\True}{\formula{A} \land \formula{B}}, just = {\TRule{\False}{\lif}[1]}
      [\sFmla{\False}{\formula{A}}, just = {\TRule{\False}{\lif}[1]}
        [\sFmla{\True}{\formula{A}}, just = {\TRule{\True}{\land}[2]}
          [\sFmla{\True}{\formula{B}}, just = {\TRule{\True}{\land}[2]}, close]
        ]
      ]
    ]
  ]
\end{oltableau}

Himpunan $\Gamma$ tidak konsisten dalam kalkulus !!{tableau} jika
terdapat suatu !!{tableau} tertutup untuk asumsi-asumsi
\[
\{\sFmla{\True}{!B_1}, \dots, \sFmla{\True}{!B_n}\}
\]
dengan syarat $!B_i \in \Gamma$ berlaku untuk setiap indeks dalam
daftar berhingga tersebut.

!!^{tableau}s diciptakan pada tahun 1950-an secara independen oleh
Evert Beth dan Jaakko Hintikka, lalu disederhanakan dan dipopulerkan
oleh Raymond Smullyan. Tableau sangat mudah digunakan karena
penyusunan !!a{tableau} merupakan prosedur yang sangat sistematis.
Karena sifat sistematis !!{tableau}s, tableau juga mudah
diimplementasikan dengan komputer. Namun, !!a{tableau} sering kali
sulit dibaca dan kaitannya dengan pembuktian kadang-kadang tidak mudah
dilihat. Pendekatan ini juga cukup umum, dan banyak logika berbeda yang
memiliki sistem !!{tableau}. !!^{tableau}s juga membantu kita
menemukan !!{structure}s yang memenuhi (himpunan) !!{sentence}s
tertentu: jika himpunan itu dapat dipenuhi, himpunan itu tidak akan
memiliki !!{tableau} tertutup, yaitu setiap !!{tableau} akan memiliki
cabang terbuka. Adapun !!{structure} yang memenuhi dapat ``dibaca'' dari
cabang terbuka, asalkan setiap aturan yang mungkin diterapkan pada
cabang itu telah diterapkan. Terdapat pula kaitan yang sangat erat
dengan kalkulus sekuen: pada dasarnya, suatu !!{tableau} tertutup merupakan
!!{derivation} kalkulus sekuen yang dipadatkan dan ditulis terbalik.

\end{document}
